\section{Specification and design}
\label{sec:method}

This section defines the experimental setting, threat model, attack variants, explanation methods, and evaluation metrics used in \Cref{sec:evaluation}.
The design aims to answer the research questions with a controlled and reproducible protocol rather than to maximise attack performance on every possible configuration.

\subsection{Dataset, model, and preprocessing}
All experiments use the CIFAR-10 test set \cite{krizhevsky2009cifar} with a fixed index range of 10{,}000 images (indices 0--9{,}999).
A fixed pretrained CIFAR-10 classifier is used throughout, specifically a ResNet-20 variant as configured in the project code.
Inputs are $32\times 32$ RGB images, and the same preprocessing and normalisation pipeline is used across all runs.

An image is considered \emph{eligible} if the clean model classifies it correctly.
Attack success rates are computed over this eligible subset rather than over the full 10k index range, so that the attack is evaluated only on images for which there is a correct clean prediction to overturn.

\subsection{Threat model}
The attacker is \emph{black-box}: model parameters and gradients are unavailable, and the attack may interact with the classifier only through forward queries that return output scores.
The perturbation budget is restricted to \emph{one pixel}.
An adversarial example therefore differs from the clean input at exactly one location $(x,y)$, while all three RGB values at that location may be changed.

This threat model is intentionally strict.
It isolates the optimisation challenge of sparse black-box perturbation and makes the resulting perturbation easy to inspect.
The cost of an attack is measured in \emph{queries}, where one query corresponds to one forward evaluation of the model on a candidate perturbed image.

\subsection{Attack variants}
All attack variants are based on Differential Evolution (DE) \cite{storn1997de}, following the one-pixel attack formulation of Su et al.\ \cite{su2019onepixel}.
A candidate solution encodes a pixel location $(x,y)$ and RGB values $(r,g,b)$.
DE evolves a population of such candidates through mutation, crossover, and selection according to the objective of the chosen attack mode.

\subsubsection{Untargeted one-pixel DE (baseline)}
The baseline attack is an untargeted DE search that attempts to induce any misclassification with a single-pixel change.
For the frozen 10k-scale run, the configuration uses a population size of 400, differential weight $F=0.5$, crossover rate $Cr=0.9$, a maximum of 100 generations, early-stopping thresholds \texttt{es\_target}=0.9 and \texttt{es\_nonadv}=0.05, and random seed 1337.

A trial is counted as successful if the final perturbed image is adversarial with respect to the eligible clean prediction.
The output CSV records, for each image, whether the attack succeeds, the number of generations used, the query count, and the best one-pixel parameters found.
This baseline provides the reference point for both research questions.

\subsubsection{Fastprior multi-target DE (fixed-budget)}
To study a guidance-style variant with a distinct cost profile, the project evaluates a fixed-budget multi-target configuration referred to as \emph{fastprior}.
For each eligible image, the attack considers target labels different from the clean prediction and runs DE under a fixed budget per target.

For the frozen 10k-scale run, the recorded configuration uses population size \texttt{pop}=192, maximum generations $=40$, $F=0.5$, $Cr=0.9$, \texttt{es\_target}$=0.9$, spatial prior controls \texttt{xy\_topk}$=32$ and \texttt{xy\_mutate\_prob}$=0.1$, and random seed 1337.
An image is counted as successful for the \emph{image-wise} attack success rate if at least one target succeeds, which is encoded by a non-zero \texttt{hit\_mask} in the output.
Because the budget is fixed, this variant has an almost constant query cost per image in the frozen results.
This makes it useful for analysing whether a stronger prior can produce more structured successful cases, even when the image-wise success rate is reduced.

\subsubsection{Responsibility-map-guided DE (RISE-guided)}
A second guided variant uses a black-box responsibility map to bias the DE search toward pixel locations predicted to be influential for the current model decision.
The responsibility map is estimated using RISE \cite{petsiuk2018rise}.
This saliency map is converted into a sampling distribution over pixel locations and used to bias spatial exploration during DE.

The purpose of this design is not to change the one-pixel threat model, but to change \emph{where} the search spends effort.
If a single-pixel attack benefits from concentrating queries on important regions, then such guidance should improve the success--cost trade-off.
If not, the guided variant may preserve or even worsen cost while leaving headline success unchanged.
This question is assessed directly in \Cref{sec:evaluation} using the frozen 10k-scale outputs under \texttt{results/final\_pack/}.

\subsection{Explanation methods}
To analyse explanation behaviour before and after attack, three explanation methods are computed for the same image index on the clean and adversarial inputs:
\begin{itemize}
  \item \textbf{Grad-CAM} \cite{selvaraju2017gradcam}: gradient-weighted localisation based on a late convolutional layer, yielding a coarse class-specific heatmap.
  \item \textbf{Integrated Gradients (IG)} \cite{sundararajan2017ig}: path-integrated attributions from a baseline input to the observed input.
  \item \textbf{RISE} \cite{petsiuk2018rise}: black-box attribution obtained by random masking and aggregation of output responses.
\end{itemize}

These methods were chosen because they cover both gradient-based and black-box explanation settings, and because they differ in granularity and stochasticity.
This makes them appropriate for testing whether explanation stability under attack is consistent across explanatory mechanisms.

\subsection{Metrics}
\subsubsection{Attack effectiveness and cost}
Attack effectiveness is measured by \emph{attack success rate} (ASR), computed over the eligible image set.
For fastprior, both image-wise success and target-wise success are relevant: image-wise success asks whether at least one target is achieved for a given image, while target-wise success counts successful targets across all target attempts.
Attack cost is measured in model queries.

The report emphasises median query counts, reported over (i) all eligible images and (ii) successful images.
This choice makes comparisons less sensitive to skewed cost distributions and allows fixed-budget and adaptive-stopping configurations to be compared in a consistent way.

\subsubsection{Explanation stability}
Explanation stability is measured on successful adversarial examples using two complementary statistics:
\begin{itemize}
  \item \textbf{IoU@10}: the intersection-over-union of the top-$k$ salient pixels before and after attack, with $k=10$.
  \item \textbf{Spearman rank correlation} $\rho$: the rank correlation between flattened attribution maps before and after attack.
\end{itemize}
IoU@10 captures whether the most salient pixels remain similar, while Spearman $\rho$ captures whether the overall attribution ordering remains similar.

\subsubsection{Explanation faithfulness via deletion AUC}
Faithfulness is assessed using deletion AUC.
Pixels are removed in descending order of saliency, and the model score for the relevant class is tracked as more of the image is deleted.
The area under this deletion curve is computed for clean and adversarial inputs.
The report presents both the clean/adversarial values and the difference $\Delta$ deletion AUC (adversarial minus clean), where a more negative change indicates a larger drop in faithfulness under attack.

\subsection{Experimental protocol and aggregation}
All runs use the same fixed 10k index range.
Eligibility is determined once on the clean model, and attack metrics are computed over the resulting eligible subset.
Explanation metrics are computed only on successful adversarial examples.
Consequently, the reported stability and faithfulness values reflect both the perturbations themselves and the composition of successful cases produced by each attack variant.

To keep the report internally consistent, all final numbers are taken from the frozen results pack under \texttt{results/final\_pack/}.
Random seeds are fixed for the frozen runs (seed 1337), and the tables and plots used in the report are generated from fixed CSV artefacts rather than transcribed manually.
